Blink amplitude varies across subjects, electrodes, and sessions, so a single global threshold becomes sensitive to the composition of epochs used to estimate amplitude statistics \citep{croft2000removal,guttmann2019new}. When blink-free and blink-heavy epochs are pooled, the mixture dilutes the amplitude distribution and may distort the global threshold, which can lower recall by making smaller or session-specific blinks less likely to exceed the decision boundary \citep{klein2013reliable,guttmann2019new}. Stage A addresses this potential bias by concentrating Stage B estimation on suspicious epochs rather than on the full clean--blink mixture, while the median/MAD estimator resists within-epoch peak contamination during threshold estimation \citep{klein2013reliable}. This mechanism is consistent with why Proposed-Med edges Proposed-Mean (pooled macro-$F_1$ 0.8667 for Proposed-Med vs 0.8534 for Proposed-Mean) without requiring a stronger claim about universal superiority.
